\subsection{GRADED TENSOR PRODUCT OF GRADED MODULES}

Let \(\Delta\) be a commutative monoid with its identity element denoted by 0, A a graded ring of type \(\Delta\) and M (resp. N) a graded right (resp. left) A-module of type \(\Delta\) . Let \((\mathbf{A}_{\lambda})\) (resp. \((\mathbf{M}_{\lambda}), (\mathbf{N}_{\lambda}))\) be the graduation of A (resp. M, N); the tensor product \(\mathbf{M} \otimes_{\mathbf{Z}} \mathbf{N}\) of the Z-modules M and N is the direct sum of the \(\mathbf{M}_{\lambda} \otimes_{\mathbf{Z}} \mathbf{N}_{\mu}\) (§ 3, no. 7, Proposition 7) and hence the latter form a bigraduation of types \(\Delta, \Delta\) on this Z-module. Consider on M \(\otimes_{\mathbf{Z}} \mathbf{N}\) the total graduation of type \(\Delta\) associated with this bigraduation (no. 1, Example 4); it consists of the sub-Z-modules \(\mathrm{P}_{\lambda} = \sum_{\mu + v = \lambda} (\mathrm{M}_{\mu} \otimes_{\mathbf{z}} \mathrm{N}_{v})\) . It is known that the Z-module M \(\otimes_{\mathbf{A}} \mathbf{N}\) is the quotient of M \(\otimes_{\mathbf{Z}} \mathbf{N}\) by the sub-Z-module Q generated by the elements \((xa) \otimes y - x \otimes (ay)\) , where \(x \in \mathbf{M}, y \in \mathbf{N}\) and \(a \in \mathbf{A}\) (§ 3, no. 1); if, for all \(\lambda \in \Delta\) , \(x_{\lambda}, y_{\lambda}, a_{\lambda}\) are the homogeneous components of degree \(\lambda\) of \(x, y, a\) respectively, clearly \((xa) \otimes y - x \otimes (zy)\) is the sum of the homogeneous elements \((x_{\lambda}a_{v}) \otimes y_{\mu} - x_{\lambda} \otimes (a_{v}y_{\mu})\) , in other words Q is a graded sub-Z-module of M \(\otimes_{\mathbf{Z}} \mathbf{N}\) (no. 3, Proposition 2) and the quotient

\[
\mathbf {M} \otimes_ {\mathbf {A}} \mathbf {N} = (\mathbf {M} \otimes_ {\mathbf {Z}} \mathbf {N}) / \mathbf {Q}
\]

therefore has canonically a graded Z-module structure of type \(\Delta\) (no. 3). Moreover (no. 3, Proposition 5), the centre C of A is a graded subring of A; the graduation which we have just defined on M \(\otimes_{\mathbf{A}} \mathbf{N}\) is compatible with its module structure over the graded ring C. For M \(\otimes_{\mathbf{Z}} \mathbf{N}\) has canonically two C-module structures, for which respectively \(c(x \otimes y) = (xc) \otimes y\) and \((x \otimes y)c = x \otimes (cy)\) for \(x \in \mathbf{M}, y \in \mathbf{N}, c \in \mathbf{C}\) ( \(\S 3\) , no. 3); if \(x \in \mathbf{M}_{\lambda}, y \in \mathbf{N}_{\mu}, c \in \mathbf{C} \cap \mathbf{A}_{\nu}\) , the two elements \(c(x \otimes y)\) and \((x \otimes y)c\) belong to (M \(\otimes_{\mathbf{Z}} \mathbf{N})_{\lambda + \mu + \nu}\) and their difference belongs to Q and hence their common image in M \(\otimes_{\mathbf{A}} \mathbf{N}\) belongs to (M \(\otimes_{\mathbf{A}} \mathbf{N})_{\lambda + \mu + \nu}\) , which establishes our assertion. When we speak of M \(\otimes_{\mathbf{A}} \mathbf{N}\) as a graded C-module, we always mean with the structure thus defined, unless otherwise mentioned. Note that (M \(\otimes_{\mathbf{A}} \mathbf{N})_{\lambda}\) can be defined as the additive group of M \(\otimes_{\mathbf{A}} \mathbf{N}\) generated by the \(x_{\mu} \otimes y_{\nu}\) , where \(x_{\mu} \in \mathbf{M}_{\mu}, y_{\nu} \in \mathbf{N}_{\nu}\) and \(\mu + \nu = \lambda\) .

Let \(\mathbf{M}'\) (resp. \(\mathbf{N}'\) ) be another graded right (resp. left) A-module and \(u: \mathbf{M} \to \mathbf{M}', v: \mathbf{N} \to \mathbf{N}'\) graded homomorphisms of respective degrees \(\alpha\) and \(\beta\) . Then it follows immediately from the above remark that \(u \otimes v\) is a graded (C-module) homomorphism of degree \(\alpha + \beta\) .

When A is commutative, a graduation (compatible with the A-module structure) is similarly defined on the tensor product of any finite number of graded A-modules; it is moreover immediate that the associativity isomorphisms such as \((\mathbf{M} \otimes \mathbf{N}) \otimes \mathbf{P} \to \mathbf{M} \otimes (\mathbf{N} \otimes \mathbf{P})\) (§ 3, no. 8, Proposition 8) are isomorphisms of graded modules.

Remark. When A has the trivial graduation (no. 1, Example 1), \((\mathbf{M} \otimes_{\mathbf{A}} \mathbf{N})_{\lambda}\) is then simply the direct sum of the sub-Y-modules \(\mathbf{M}_{\mu} \otimes_{\mathbf{A}} \mathbf{N}_{\nu}\) of \(\mathbf{M} \otimes_{\mathbf{A}} \mathbf{N}\) such that \(\mu + \nu = \lambda\) .

Let M (resp. N) be a graded right (resp. left) A-module of type \(\Delta\) , P a graded Z-module of type \(\Delta\) and let f be a Z-bilinear mapping of \(M \times N\) into P satisfying condition (1) of §3, no. 1, and such moreover that

\[
f (x _ {\lambda}, y _ {\mu}) \in \mathrm{P} _ {\lambda + \mu} \quad \text { for } x _ {\lambda} \in \mathrm{M} _ {\lambda}, y _ {\mu} \in \mathrm{N} _ {\mu}, \lambda , \mu \text { in } \Delta .
\]

Then \(f(x, y) = g(x \otimes y)\) , where \(g: \mathbf{M} \otimes_{\mathbf{A}} \mathbf{N} \to \mathbf{P}\) is a \(\mathbf{Z}\) -linear mapping (§ 3, no. 1, Proposition 1) and it follows from the above condition that \(g\) is a graded \(\mathbf{Z}\) -module homomorphism of degree 0.

Let B be another graded ring of type \(\Delta\) and \(\rho: A \to B\) a homomorphism of graded rings (no. 2); then \(\rho_*(B_d)\) is a graded right A-module of type \(\Delta\) . If E is a graded left A-module of type \(\Delta\) and \(\rho^*(B_d) \otimes_A E\) is given the graded Z-module structure of type \(\Delta\) defined above, the canonical left B-module structure ( \(\S\) 5, no. 1) is compatible with the graduation of

\[
\mathbf {E} _ {(\mathbf {B})} = \rho^ {*} (\mathbf {E}) = \rho_ {*} (\mathbf {B} _ {d}) \otimes_ {\mathbf {A}} \mathbf {E}.
\]

The graded B-module thus obtained is said to be obtained by extending the ring of scalars to B by means of \(\rho\) and when we speak of \(\mathrm{E}_{(\mathrm{B})}\) or \(\rho^{*}(\mathrm{E})\) as a graded B-module, we always mean this structure, unless otherwise mentioned.

